\section{Motivation}
\label{sec:intro:landscape}

% Robot manipulation systems now operate in environments that were long considered beyond the reach of automation.
Robotic manipulation has moved from structured factory floors into increasingly complex, variable, and human-centered environments.
Robots routinely sort and transport goods through warehouses, assemble products on factory lines, assist in surgical and healthcare settings, prepare food, and execute tasks in hazardous and remote environments.
At the same time, the field has expanded beyond rigid industrial automation toward increasingly general-purpose manipulation capabilities: robots are now expected not only to repeat fixed motions reliably, but also to adapt to novel objects, changing environments, and unforeseen task conditions.
Meeting this expectation requires robots to generate actions that are not only task-directed, but also executable under the kinematic and dynamic (define) constraints imposed by the robot-world system.
For instance, a manipulator asked to slide a heavy object across a cluttered table must respect reachability, contact geometry, friction, actuation limits, and the inertial response of the object once contact is broken.
These constraints are therefore not peripheral to manipulation; they determine which states are reachable, which transitions are physically realizable, and which regions of the motion space are worth exploring at all.
The technological progress enabling modern manipulation systems has followed two largely separate but equally influential trajectories.
(citations)

\header{Engineered pipelines at scale}
Industrial manipulation systems are typically organized as engineered pipelines: perception localizes objects, planning produces motions, and control executes them.
These systems have enabled reliable deployment across warehouses, factories, surgical and assistive settings, chemical and pharmaceutical processing, food service, and extreme environments.
Their strength lies in modularity, debuggability, and the accumulation of engineering effort over repeated deployment cycles.

\header{General-purpose learned policies}
A parallel line of work seeks to reduce this engineering burden through general-purpose learned policies trained on large and diverse datasets.
These systems can generalize across objects, scenes, and modest task variation without hand-designing a new pipeline for every setting.

\header{Prehensile manipulation and the geometric paradigm}
Underlying both paradigms is a shared assumption that has shaped manipulation research for decades: that manipulation is fundamentally a geometric problem.
The dominant objects of reasoning are collision clearance, grasp pose, reachability, and singularity avoidance---broadly, kinematic constraints on robot configuration.
This assumption is not arbitrary.
Grasping has attracted disproportionate research attention \cite{billard2019trends,bohg2013data,mahler2017dex} precisely because form-closure reduces object state to a geometry problem: once an object is secured, its state is predictable and planning reduces to obstacle avoidance \cite{mason1986mechanics,bullock2011classifying}.
Recent years have brought substantial improvements to geometric planning infrastructure---GPU-parallelized collision checking and forward kinematics \cite{sundaralingam2023curobo,thomason2024motions}, large-scale motion datasets \cite{o2024open}, and high-fidelity simulation environments \cite{coumans2021}---but the underlying planning abstraction has remained largely geometric.

\header{Behaviors beyond geometric planning}
The cost of this assumption becomes visible at the boundary of kinematic behaviors.
A broad class of manipulation tasks is not adequately captured by geometry alone: tasks where joint torques, object inertia, contact impulses, and fluid dynamics jointly determine the outcome.
Non-prehensile rearrangement through impulsive contact \cite{pasricha2022pokerrt,mason2018toward}, transport of liquid-filled containers \cite{abderezaei2024clutter}, high-speed inertial transport \cite{ichnowski2022gomp,zeng2020tossingbot}, operation near or beyond rated payload capacity, and adaptation to mid-task mechanical failures \cite{briscoe2024exploring} all fall into this category.
These behaviors require robot control to reason over contact, to model relevant physics, and to search over a higher-dimensional space than kinematic planning allows.

\header{Expanding manipulation capability}
Accessing the full range of manipulation capability afforded by modern robot embodiments---beyond the kinematic behaviors that dominate current systems---requires explicitly modeling and integrating dynamic constraints into the motion generation process.
Tasks involving complex object interactions, unstructured environments, and high-speed or high-load operations cannot be adequately planned or controlled using geometry alone.
Expanding manipulation capability requires moving beyond kinematic behaviors to model and plan under dynamic constraints, enabling robots to reason about their available embodiment and complex object interactions.

The constraints governing robot manipulation decompose into two families.
\textit{Kinematic constraints}---joint angle limits, self-collision avoidance, environmental obstacle clearance, singularity avoidance, and reachability---define the geometry of feasible motion.
\textit{Dynamic constraints} govern the physics of motion and arise from two sources: \textit{higher-order robot constraints}---bounds on velocity, acceleration, jerk, and torque that arise from the robot's own kinematics and actuator limits \cite{pham2018new,berscheid2021jerk}---and \textit{task-induced dynamic constraints} that arise from the physics of robot-object interaction, including contact mechanics, inertial coupling, end-effector acceleration limits, and motion-induced phenomena such as sloshing.
